\section*{Bab \ref{ch:b2:aromatic-substitution} --- Aromatisitas dan Substitusi Aromatik Elektrofilik}

\begin{solution}{exo:b2:aromatic-substitution:1}
\ce{HNO3 + 2H2SO4 -> NO2+ + H3O+ + 2HSO4-}: asam nitrat terprotonasi, lalu melepaskan air.
\end{solution}

\begin{solution}{exo:b2:aromatic-substitution:2}
Toluena menghasilkan 2-bromotoluena dan 4-bromotoluena: metil mengarahkan ke
posisi orto dan para. Nitrobenzena menghasilkan 3-bromonitrobenzena, lebih
lambat: gugus nitro mengarahkan ke meta.
\end{solution}

\begin{solution}{exo:b2:aromatic-substitution:3}
\ce{-CH3}: pengaktif, o/p. \ce{-OH}: pengaktif kuat, o/p. \ce{-Cl}: pendeaktif, o/p.
\ce{-NO2}: pendeaktif kuat, meta. \ce{-COCH3}: pendeaktif, meta. \ce{-NHCOCH3}:
pengaktif (sedang), o/p.
\end{solution}

\begin{solution}{exo:b2:aromatic-substitution:4}
Asetofenon (1-feniletan-1-on), \ce{C6H5COCH3}, melalui \omterm{def:b2:aromatic-substitution:friedel-crafts}{ion asilium} \ce{CH3CO+}.
\end{solution}

\begin{solution}{exo:b2:aromatic-substitution:5}
Karbokation primer (atau kompleksnya) mengalami penataan ulang melalui
pergeseran hidrida menjadi kation isopropil sekunder, yang mengalkilasi
cincin. Propilbenzena: \omterm{def:b2:aromatic-substitution:friedel-crafts}{asilasi Friedel--Crafts} dengan propanoil klorida
(tanpa penataan ulang), lalu reduksi \ce{C=O} keton menjadi \ce{CH2}.
\end{solution}

\begin{solution}{exo:b2:aromatic-substitution:6}
Gugus \ce{-OH} adalah donor mesomerik yang kuat: cincinnya begitu teraktivasi
sehingga brom molekular sudah merupakan elektrofil yang memadai, dan ketiga
posisi orto dan para tersubstitusi semuanya.
\end{solution}

\begin{solution}{exo:b2:aromatic-substitution:7}
3-Bromonitrobenzena: nitrasi, lalu brominasi (nitro mengarahkan ke meta).
4-Bromonitrobenzena: brominasi, lalu nitrasi (brom mengarahkan ke
orto/para), dan pisahkan isomer para dari isomer orto.
\end{solution}

\begin{solution}{exo:b2:aromatic-substitution:8}
Campuran penitrasi mengoksidasi anilina yang sangat kaya elektron, dan
memprotonasi nitrogennya: \ce{-NH3+} adalah \omterm{def:b2:aromatic-substitution:directing}{gugus pendeaktif} \omterm{def:b2:aromatic-substitution:directing}{pengarah meta}.
Dalam asetanilida, pasangan elektron bebas nitrogen dipakai bersama dengan
karbonil: karena kurang basa dan kurang mengaktifkan, nitrogen itu tidak
terprotonasi atau teroksidasi, dan tetap mengarahkan ke orto/para, terutama
para karena alasan sterik.
\end{solution}

\begin{solution}{exo:b2:aromatic-substitution:9}
Sulfonasi toluena (terutama para, karena gugus yang meruah menghindari
orto); brominasi: brom masuk pada posisi orto terhadap metil (posisi para
terblokir, dan kedua gugus menyukai posisi itu); lepaskan gugus \ce{SO3H}
dengan pemanasan dalam asam encer.
\end{solution}

\begin{solution}{exo:b2:aromatic-substitution:10}
Kedua gugus mengarahkan ke orto/para; posisi para masing-masing ditempati
oleh gugus yang lain. Metoksi, pengaktif yang lebih kuat, menang: nitrasi
terjadi pada posisi orto terhadap \ce{-OCH3} (posisi 2), dan menghasilkan
4-metil-2-nitroanisol.
\end{solution}

\begin{solution}{exo:b2:aromatic-substitution:11}
Setiap gugus nitro sangat mendeaktifkan cincin: nitrasi kedua memerlukan
kondisi yang lebih keras daripada yang pertama, dan yang ketiga (pada
cincin yang membawa dua gugus nitro dan hanya metil yang mengaktifkan
secara lemah) jauh lebih keras lagi.
\end{solution}

\begin{solution}{exo:b2:aromatic-substitution:12}
Asam 4-nitrobenzoat: nitrasi toluena, pisahkan isomer para, oksidasi
metilnya menjadi \ce{-COOH} (permanganat panas). Asam 3-nitrobenzoat:
oksidasi toluena menjadi asam benzoat lebih dahulu, lalu nitrasi
(\ce{-COOH} mengarahkan ke meta).
\end{solution}

\begin{solution}{pb:b2:aromatic-substitution:1}
\textbf{1.} \ce{HNO3 + 2H2SO4 -> NO2+ + H3O+ + 2HSO4-}.
\textbf{2.} \ce{O=N+=O}: linear, isoelektronik dengan \ce{CO2}.
\textbf{3.} Nitrogen berikatan dengan karbon cincin yang para terhadap
\ce{CH3}, yang menjadi tetrahedral (membawa H dan \ce{NO2}).
\textbf{4.} Muatan positif pada kedua karbon yang orto terhadap karbon yang
diserang dan pada karbon yang para terhadapnya, yang membawa metil.
\textbf{5.} Tahap pertama, pembentukan \omterm{def:b2:aromatic-substitution:seAr}{zat antara Wheland}.
\textbf{6.} Basa dalam medium: \ce{HSO4-} atau air.
\textbf{7.} Dua orto, dua meta, satu para.
\textbf{8.} 2 : 2 : 1, yaitu 40~\% orto, 40~\% meta, 20~\% para.
\textbf{9.} Pada serangan meta, muatan positif tidak pernah berada pada
karbon yang membawa metil; tidak ada struktur yang dibantu oleh metil.
\textbf{10.} Satu struktur resonansi mempunyai muatan pada karbon yang
membawa \ce{CH3}: karbokation tersier, yang distabilkan oleh efek induktif
dan \omterm{def:b2:fragment-orbitals:hyperconjugation}{hiperkonjugasi} metil.
\textbf{11.} Orto dan para disukai (96~\% dibandingkan dengan 60~\% secara
acak); meta hampir tidak ada.
\textbf{12.} Per posisi: orto $58/2 = 29~\%$, para 38~\%: para adalah posisi yang lebih reaktif,
karena posisi-posisi orto sedikit terhalang oleh metil.
\textbf{13.} Lebih cepat: metil mengaktifkan cincin.
\textbf{14.} 4-Nitrotoluena (meleleh pada \qty{52}{\degreeCelsius}).
\textbf{15.} Dinginkan campuran mentahnya: isomer para mengkristal dan
disaring; cairannya mengandung isomer orto dan meta.
\textbf{16.} Molekul yang simetris lebih rapat tersusun dalam kristal: lebih
banyak interaksi per volume, titik leleh lebih tinggi.
\textbf{17.} Dengan \omterm{def:b2:liquid-vapour-diagrams:fractional-distillation}{distilasi fraksional} pada tekanan rendah, atau dengan
kromatografi; atau dengan pendinginan lebih lanjut (isomer meta meleleh pada
\qty{16}{\degreeCelsius}).
\textbf{18.} Titik didih isomer-isomer itu berdekatan (rumus sama,
kepolaran mirip).
\textbf{19.} \ce{C7H8}: \qty{92.14}{g/mol}; \ce{C7H7NO2}: \qty{137.14}{g/mol}.
\textbf{20.} $100/92.14 = \qty{1.085}{mol}$ toluena; $0.90 \times 1.085 = \qty{0.977}{mol}$
nitrotoluena.
\textbf{21.} Total $0.977 \times 137.14 = \qty{134.0}{g}$: orto \qty{77.7}{g}, meta
\qty{5.4}{g}, para \qty{50.9}{g}.
\textbf{22.} Nitrasi kedua (dinitrotoluena), dan oksidasi gugus metil.
\textbf{23.} Oksidasi metilnya menjadi \ce{-COOH} dengan permanganat basa
panas, lalu asamkan.
\textbf{24.} Percobaan itu menghasilkan $\boldsymbol{\approx \qty{51}{g}}$ 4-nitrotoluena.
\end{solution}
